\section{Критерии качества}

\begin{frame}[c]{Как понять, что модель хороша}
  \small
  \zebra{1}
  \begin{tabular}{>{\raggedright\arraybackslash}p{0.191\textwidth}>{\raggedright\arraybackslash}p{0.688\textwidth}}
    \textcolor{acMid}{\textbf{Бизнес}} & Сколько денег и довольных клиентов получает компания. Зависит от всего сразу, измерить трудно.\\
    \textcolor{acMid}{\textbf{Онлайн}} & Что видно на работающем сервисе: например, сколько времени люди сидят в приложении.\\
    \textcolor{acMid}{\textbf{Асессоры}} & Живые люди смотрят на ответы модели и ставят оценки. Нужно, когда качество --- дело вкуса.\\
    \textcolor{acMid}{\textbf{Офлайн}} & Считаем по старым данным до запуска: насколько ответы модели отличаются от правильных.\\
  \end{tabular}
\end{frame}

\begin{frame}[c]{Офлайн и онлайн метрики}
  \small
  \zebra{2}\small
  \begin{tabular}{>{\raggedright\arraybackslash}p{0.115\textwidth}>{\raggedright\arraybackslash}p{0.362\textwidth}>{\raggedright\arraybackslash}p{0.382\textwidth}}
    \rowcolor{acDark} & \textcolor{white}{\textbf{Офлайн}} & \textcolor{white}{\textbf{Онлайн}}\\
    Когда & До запуска модели & После деплоя, на работающей системе\\
    На чём & Исторические данные, тестовая выборка, оценки асессоров & Реальное поведение пользователей и магазинов\\
    Примеры & accuracy, MAE, MSE, precision, recall & Медианная длина сессии, число бананов на полках в конце дня\\
    Плюсы & Быстро, дёшево, без риска для пользователей & Ближе к бизнес-метрикам\\
    Минусы & Хороший результат не гарантирует рост бизнес-метрик & Нужны деплой, A/B-тест и риск для пользователей\\
  \end{tabular}
\end{frame}

\begin{frame}[c]{Путь модели от офлайна до бизнеса}
  \lead{Как всё связано.}

  \medskip
  \begin{columns}[T]
    \begin{column}{0.23\textwidth}
      \begin{block}{Офлайн}
        Метрики на тесте. Быстро отсеиваем слабые модели.
      \end{block}
    \end{column}
    \begin{column}{0.23\textwidth}
      \begin{block}{Деплой}
        Выкатка в продакшен. Нужна возможность отката.
      \end{block}
    \end{column}
    \begin{column}{0.23\textwidth}
      \begin{block}{Онлайн и A/B}
        Сравниваем с прошлой версией на случайных группах.
      \end{block}
    \end{column}
    \begin{column}{0.23\textwidth}
      \begin{block}{Бизнес}
        Доход и лояльность --- итоговая цель.
      \end{block}
    \end{column}
  \end{columns}

  \medskip
  \textbf{Офлайн-метрика --- прокси онлайн-метрики, онлайн --- прокси бизнеса.
  Хорошее число на тесте не гарантирует эффект в продакшене.}
\end{frame}

\begin{frame}[c]{Офлайн или онлайн?}
  \lead{Голосуем.}
  \begin{numrows}
        1 & Доля верно поставленных диагнозов на отложенной выборке\\
        2 & Среднее число бананов на полках в конце дня\\
        3 & Средний модуль разницы между предсказанной и истинной ценой квартиры\\
        4 & Медианная длина сессии в онлайн-игре\\
        5 & Оценка качества перевода асессорами\\
      \end{numrows}\par
\end{frame}

\begin{frame}[c]{Ответы: вид метрики}
  \small
  \zebra{2}
  \begin{tabular}{>{\raggedright\arraybackslash}p{0.324\textwidth}>{\raggedright\arraybackslash}p{0.553\textwidth}}
    \rowcolor{acDark} \textcolor{white}{\textbf{Метрика}} & \textcolor{white}{\textbf{Вид}}\\
    1. Доля верных диагнозов & Офлайн: считаем по историческим данным\\
    2. Бананы на полках & Онлайн: нужна работающая система\\
    3. Ошибка в цене квартиры & Офлайн: сравниваем с известными ценами\\
    4. Длина сессии & Онлайн: считается по реальным пользователям\\
    5. Оценка асессоров & Офлайн: оценивают до выпуска в продакшен\\
  \end{tabular}
\end{frame}

\begin{frame}[c]{Какую метрику выбрать?}
  \lead{Вопрос на подумать.}
  \small Модель отвечает «да» или «нет». Ошибиться можно двумя способами: сказать «да» зря
  (ложная тревога) или сказать «нет» зря (пропуск).

  \smallskip
  \begin{columns}[T]
    \begin{column}{0.31\textwidth}
      \begin{block}{Обычная поликлиника}
        Нужно, чтобы диагнозы в целом чаще были верными.
      \end{block}
    \end{column}
    \begin{column}{0.31\textwidth}
      \begin{block}{Проходная завода}
        Камера решает, сотрудник ли перед ней. Пустить чужого --- беда, а лишний раз не пустить своего --- мелочь.
      \end{block}
    \end{column}
    \begin{column}{0.31\textwidth}
      \begin{block}{Опасная болезнь}
        Пропущенный больной заразит других. Ложная тревога куда дешевле.
      \end{block}
    \end{column}
  \end{columns}

  \smallskip
  \textbf{Какую метрику выбрать в каждом случае?}
\end{frame}

\begin{frame}[c]{Ответ: разные нужды --- разные метрики}
  \small
  \zebra{1}
  \begin{tabular}{>{\raggedright\arraybackslash}p{0.191\textwidth}>{\raggedright\arraybackslash}p{0.688\textwidth}}
    \textcolor{acMid}{\textbf{Accuracy}}\newline{\scriptsize\color{acGray}обычная поликлиника} &
    \textbf{верные ответы $\div$ все ответы}\newline
    Какая доля ответов оказалась правильной. Подходит, когда обе ошибки одинаково неприятны.\\
    \textcolor{acMid}{\textbf{Precision}}\newline{\scriptsize\color{acGray}проходная завода} &
    \textbf{верно найденные $\div$ всё, что модель назвала «да»}\newline
    Можно ли верить модели, когда она говорит «да». Важна, когда ложное «да» дорого стоит: нельзя пустить чужого.\\
    \textcolor{acMid}{\textbf{Recall}}\newline{\scriptsize\color{acGray}опасная болезнь} &
    \textbf{верно найденные $\div$ все, кого нужно было найти}\newline
    Сколько из настоящих «да» модель нашла. Важна, когда нельзя пропустить: больной заразит других.\\
  \end{tabular}
\end{frame}

\begin{frame}[c]{Другие метрики}
  \lead{Вопрос на подумать.} Какую метрику выбрать здесь?

  \smallskip
  \begin{columns}[T]
    \begin{column}{0.31\textwidth}
      \begin{block}{Видеозвонок}
        Приложение «надевает» на участника виртуальный костюм рыцаря. Костюм должен ровно лечь на фигуру.
      \end{block}
    \end{column}
    \begin{column}{0.31\textwidth}
      \begin{block}{Безопасность}
        Охрана важного объекта: нужно заметить любого, кто вошёл в кадр.
      \end{block}
    \end{column}
    \begin{column}{0.31\textwidth}
      \begin{block}{Пандемия}
        Контроль, не превышает ли число людей в помещении порог.
      \end{block}
    \end{column}
  \end{columns}
\end{frame}

\begin{frame}[c]{Ответ: зависит от цели}
  \begin{columns}[T]
    \begin{column}{0.31\textwidth}
      \begin{block}{Точность положения}
        Среднеквадратичное отклонение координат опорных точек тела от истинных: костюм должен сесть ровно.
      \end{block}
    \end{column}
    \begin{column}{0.31\textwidth}
      \begin{block}{Полнота (recall)}
        На какой доле кадров с людьми модель их заметила. Точность рамки не так важна.
      \end{block}
    \end{column}
    \begin{column}{0.31\textwidth}
      \begin{block}{Число людей}
        Среднеквадратичное отклонение числа найденных людей от истинного.
      \end{block}
    \end{column}
  \end{columns}
\end{frame}

\begin{frame}[c]{Не всё сводится к метрикам}
  \begin{columns}[T]
    \begin{column}{0.31\textwidth}
      \begin{block}{Скорость}
        Предсказания в режиме реального времени. Требование и к модели, и к серверам.
      \end{block}
    \end{column}
    \begin{column}{0.31\textwidth}
      \begin{block}{Компактность}
        Модель должна помещаться на телефон или другое устройство.
      \end{block}
    \end{column}
    \begin{column}{0.31\textwidth}
      \begin{block}{Интерпретируемость}
        Можно объяснить решение: дадут ли кредит, согласуют ли лечение.
      \end{block}
    \end{column}
  \end{columns}
\end{frame}
